% [automatingbook](hook://file/OLBvzndSz?p=X1NjcmF0Y2hQYWQvSmF2YVNjcmlwdA==&n=automatingbook)

\begin{frame}[fragile]{module.exports}

\begin{itemize}
	\item A module can have functions (lines 1 -- 3).
	\item The function can be exported (line 4) using `modules.exports'. 
\end{itemize}

\begin{Code}{b.js}%
\begin{lstlisting}[language=javascript, firstnumber=1, label=glabels, xleftmargin=0pt, basicstyle=\scriptsize]% 

const js = () => {
  console.log('JavaScript');
}
module.exports = js;
\end{lstlisting}%   
\end{Code}%

\end{frame}

\begin{frame}[fragile]{require}

\begin{itemize}
	\item To use the module's functions, we use `require' (line 1). 
	\item The reference name (b) is used to call the method (line 3). 
\end{itemize}

\begin{Code}{a.js}%
\begin{lstlisting}[language=javascript, firstnumber=1, label=glabels, xleftmargin=0pt, basicstyle=\scriptsize]% 

const b = require('./b.js');
console.log('From a.js: running code in the file b.js'); 
b();
\end{lstlisting}%   
\end{Code}%

\begin{Command}{node a.js}%
\begin{lstlisting}[firstnumber=1, label=glabels, xleftmargin=0pt, basicstyle=\scriptsize]% 

> node a.js
From a.js: running code in the file b.js
JavaScript
\end{lstlisting}%   
\end{Command}%
\end{frame}

\begin{frame}[fragile]{exports/require}

\begin{itemize}
	\item `exports' can be interpreted as a dynamic object; we can add functions to the object (lines 3 -- 4). 
\end{itemize}

\begin{Code}{d.js}%
\begin{lstlisting}[firstnumber=1, label=glabels, xleftmargin=0pt, basicstyle=\scriptsize]% 

const add = (x, y) => x + y
const sub = (x, y) => x - y

exports.add = add
exports.sub = sub
\end{lstlisting}%   
\end{Code}%

\begin{itemize}
	\item To use the exported functions, we can use `require' to access the functions. 
\end{itemize}

\begin{Code}{c.js}%
\begin{lstlisting}[firstnumber=1, label=glabels, xleftmargin=0pt, basicstyle=\scriptsize]% 

const c = require('./d');
console.log(c.add(10, 20))
console.log(c.sub(10, 20))
\end{lstlisting}%   
\end{Code}%
\end{frame}

\begin{frame}[fragile]{}

\begin{itemize}
	\item When we use a block \{ ... \} with  `require,' we can use add/sub to reference the module's add/sub functions. 
\end{itemize}

\begin{Code}{}%
\begin{lstlisting}[firstnumber=1, label=glabels, xleftmargin=0pt, basicstyle=\scriptsize]% 

const {add, sub} = require('./d');
console.log(add(10, 20))
console.log(sub(10, 20))
\end{lstlisting}%   
\end{Code}%

\begin{Command}{node c.js}%
\begin{lstlisting}[firstnumber=1, label=glabels, xleftmargin=0pt, basicstyle=\scriptsize]% 

> node c.js
30
-10
\end{lstlisting}%   
\end{Command}%

\end{frame}